% Part: sets-functions-relations
% Chapter: arithmetization
% Section: checking-details
%
\documentclass[../../../include/open-logic-section]{subfiles}


\begin{document}
	
\olfileid{sfr}{arith}{check}
\olsection{क्रमित वलये आणि क्षेत्रे}

या प्रकरणभर काही व्याख्या ``जशा वागायला हव्यात'' तशा वागतात, असे
आपण म्हटले. या तांत्रिक परिशिष्टात त्याचा अर्थ स्पष्ट करू आणि त्या
व्याख्या ``योग्य रीतीने'' वागतात हे कसे दाखवायचे याचा आराखडा देऊ.

\olref[int]{sec} मध्ये आपण $\Int$ वरील बेरीज आणि गुणाकार व्याख्यायित
केले. व्याख्येप्रमाणे या क्रिया $\Int$ ला आपण ``इच्छित'' असलेली संरचना
देतात, हे दाखवायचे आहे. विशेषतः, येथे अभिप्रेत संरचना क्रमनिरपेक्ष
वलयाची आहे.

\begin{defn}
	\emph{क्रमनिरपेक्ष वलय} म्हणजे विशिष्ट $0$ व $1$ घटक आणि $+$ व $\times$
	क्रिया असलेला असा संच $S$, जो पुढील आठ सूत्रे पूर्ण करतो:
	\begin{align*}
		\emph{सहचारिता}&&a + (b+ c) & = (a + b) + c \\
		&& (a \times b) \times c & = a \times (b\times c)\\
		\emph{क्रमनिरपेक्षता}&&a + b &= b+ a  \\
		&&  a \times b&= b\times a\\
		\emph{अविकारक}&&a + 0 &= a \\
		&& a \times 1 &= a\\
		\emph{बेरीज व्यस्त}&&(\exists b\in S)0&=a + b\\
		\emph{वितरणता}&&a \times (b+ c ) &= (a \times b) + (a \times c)
	\end{align*}
	येथे ही सर्व सूत्रे $S$ पुरते मर्यादित सार्वत्रिक परिमाणकांनी बद्ध आहेत,
	हे अध्याहृत आहे. तसेच येथील $0$ आणि~$1$ हे घटक याच नावाच्या नैसर्गिक
	संख्या असतीलच असे नाही.
\end{defn}

म्हणून पूर्णांक क्रमनिरपेक्ष वलय बनवतात हे तपासण्यासाठी या आठ अटी पूर्ण
होतात एवढेच तपासायचे आहे. यांपैकी कोणतीही अट सिद्ध करणे {अवघड} नाही,
पण ते थोडे कष्टाचे आहे. उदाहरणार्थ, बेरजेच्या बाबतीत
\emph{सहचारिता} अशी सिद्ध करता येते:

\begin{proof}
$i, j, k \in \Int$ निश्चित करा. मग $i = \equivrep{a_1, b_1}{}$ आणि $j =
\equivrep{a_2,b_2}{}$ आणि $k = \equivrep{a_3, b_3}{}$ असे होतील अशा
$a_1, b_1, a_2, b_2, a_3, b_3 \in \Nat$ संख्या आहेत. (वाचनीयतेसाठी
``$\equivrep{x, y}{\Intequiv}$'' ऐवजी ``$\equivrep{x, y}{}$'' असे लिहितो;
या संपूर्ण विभागात आपण हेच करू.) आता:
\begin{align*}
	i + (j + k) &= \equivrep{a_1, b_1}{}+(\equivrep{a_2, b_2}{} + \equivrep{a_3, b_3}{}) \\
	&= \equivrep{a_1,  b_1}{} + \equivrep{a_2+a_3, b_2+b_3}{}\\
	&= \equivrep{a_1 + (a_2 + a_3), b_1 + (b_2 + b_3)}{}\\
	&= \equivrep{(a_1 + a_2) + a_3, (b_1 + b_2) + b_3}{}\\
	&= \equivrep{a_1 + a_2, b_1 + b_2}{} + \equivrep{a_3, b_3}{}\\
	&= (\equivrep{a_1, b_1}{} + \equivrep{a_2, b_2}{}) + \equivrep{a_3, b_3}{}\\
	&= (i+j) + k
\end{align*}
येथे $\Nat$ वरील बेरजेच्या वर्तनाचा आपण मुक्तपणे उपयोग केला आहे.
\end{proof}

त्याचप्रमाणे, \emph{बेरीज व्यस्त} असा सिद्ध करता येतो:

\begin{proof}
$i \in \Int$ निश्चित करा; म्हणजे काही $a,b \in \Nat$ साठी $i =
\equivrep{a,b}{}$. $j = \equivrep{b,a}{} \in \Int$ असे ठरवू. नैसर्गिक
संख्यांच्या वर्तनाचा उपयोग केल्यास $(a+b) + 0 = 0 + (a+b)$; म्हणून
व्याख्येनुसार $\tuple{a+b, b+a} \sim_\Int \tuple{0,0}$, आणि त्यामुळे
$\equivrep{a+b, b+a}{} = \equivrep{0, 0}{} = 0_\Int$. आता $i + j =
\equivrep{a,b}{}+\equivrep{b,a}{}=\equivrep{a+b, b+a}{}= \equivrep{0,
0}{} = 0_\Int$.
\end{proof}

आणि \emph{वितरणतेची} सिद्धता अशी आहे:

\begin{proof}
वरीलप्रमाणे, $i = \equivrep{a_1, b_1}{}$ आणि $j =
\equivrep{a_2,b_2}{}$ आणि $k = \equivrep{a_3, b_3}{}$ निश्चित करा. आता:
\begin{align*}
	i \times (j + k)
	&= \equivrep{a_1, b_1}{} \times (\equivrep{a_2,b_2}{} + \equivrep{a_3, b_3}{})\\
	&= \equivrep{a_1, b_1}{} \times \equivrep{a_2 + a_3,b_2+b_3}{}\\
	&= \equivrep{a_1  (a_2 + a_3) + b_1  (b_2+b_3), a_1  (b_2 + b_3) + b_1 (a_2 + a_3)}{}\\
	&= \equivrep{a_1 a_2 + a_1a_3 + b_1 b_2+b_1b_3, a_1 b_2 + a_1b_3 + a_2b_1 + a_3b_1}{}\\		
%		&= \equivrep{a_1 a_2 + b_1b_2 + a_1a_3 + b_1 b_2+b_1b_3, a_1 b_2 + a_1b_3 + a_2b_1 + a_3b_1}{}\\		
	&= \equivrep{a_1a_2 + b_1b_2, a_1b_2 + a_2b_1}{} + \equivrep{a_1a_3 + b_1b_3, a_1b_3 + a_3b_1}{}\\
	&= (\equivrep{a_1, b_1}{} \times \equivrep{a_2,b_2}{}) + (\equivrep{a_1, b_1}{} \times  \equivrep{a_3, b_3}{})\\
	&= (i \times j) + (i \times k)
\end{align*}
\end{proof}

उरलेल्या पाच अटी सिद्ध करण्याचे काम सरावासाठी सोडतो. ते पूर्ण केल्यावर
$\Int$ हे क्रमनिरपेक्ष वलय आहे, म्हणजे व्याख्यायित केलेली बेरीज आणि
गुणाकार अपेक्षेप्रमाणे वागतात, हे आपण दाखवलेले असेल.

\begin{prob}
$\Int$ हे क्रमनिरपेक्ष वलय आहे हे सिद्ध करा.
\end{prob}

पण आपले काम अजून संपलेले नाही. $\Int$ वरील बेरीज आणि गुणाकार यांच्याबरोबरच
आपण $\leq$ हा क्रमसंबंध व्याख्यायित केला होता; तो अपेक्षेप्रमाणे वागतो हेही
तपासले पाहिजे. अधिक नेमकेपणाने, $\Int$ हे \emph{क्रमित} वलय आहे हे दाखवले
पाहिजे.\footnote{\olref[sfr][rel][ord]{def:linearorder} मधील व्याख्या आठवा:
	पूर्ण क्रम म्हणजे परावर्ती, संक्रमक, प्रतिसममित आणि संयोजित संबंध.
	क्रमसंबंधांच्या संदर्भात संयोजिततेला कधीकधी \emph{त्रिपर्याय} म्हणतात,
	कारण कोणत्याही $a$ आणि $b$ साठी $a \leq b \lor a
	= b \lor a \geq b$.}

\begin{defn}
\emph{क्रमित वलय} म्हणजे $\leq$ या पूर्ण क्रमसंबंधानेही सज्ज असलेले
असे क्रमनिरपेक्ष वलय की:
\begin{align*}
	a \leq b &\lif a + c \leq b + c\\
	(a \leq b \land 0 \leq c) &\lif a \times c \leq b \times c
\end{align*}
\end{defn}

\begin{prob}
$\Int$ हे क्रमित वलय आहे हे सिद्ध करा.
\end{prob}

पूर्वीप्रमाणेच, रचलेला $\Int$ हा क्रमित वलय आहे हे दाखवणे कष्टाचे पण
नित्याचे काम आहे. ते आम्ही तुमच्यासाठी सोडतो.

यामुळे पूर्णांकांचे काम संपले. आता परिमेय संख्यांबद्दल अगदी तशाच गोष्टी
दाखवायच्या आहेत. विशेषतः, $+$, $\times$ आणि $\leq$ यांच्या आपल्या
व्याख्यांनुसार परिमेय संख्या क्रमित \emph{क्षेत्र} बनवतात हे दाखवायचे आहे:
\begin{defn}\ollabel{orderedfield}
\emph{क्रमित क्षेत्र} म्हणजे पुढील अटही पूर्ण करणारे क्रमित वलय:
\begin{align*}
	\emph{गुणाकार व्यस्त}& & (\forall a \in S \setminus \{0\})(\exists b \in S) a\times b& = 1
\end{align*}
\end{defn}

$\Int$ हे क्रमित वलय आहे हे दाखवल्यानंतर $\Rat$ हे क्रमित क्षेत्र आहे
हे दाखवणे सोपे पण कष्टाचे आहे.

\begin{prob}
$\Rat$ हे क्रमित क्षेत्र आहे हे सिद्ध करा.
\end{prob}

पूर्णांक आणि परिमेय संख्या हाताळल्यानंतर केवळ वास्तव संख्या उरतात.
विशेषतः, $\Real$ हे \emph{संपूर्ण} क्रमित क्षेत्र, म्हणजे संपूर्णता गुणधर्म
असलेले क्रमित क्षेत्र आहे, हे दाखवले पाहिजे. \olref[cuts]{realcompleteness}
मध्ये $\Real$ ला संपूर्णता गुणधर्म आहे हे सिद्ध केले. मात्र $\Real$ हे
क्रमित क्षेत्र आहे याची (कंटाळवाणी) तपासणी अजून करायची आहे.
%READERNOTE{MRARITH-007: गोठवलेल्या इंग्रजी वाक्यात “run through the (tedious) of checking” येथे आवश्यक नाम गाळले आहे. मराठीत अभिप्रेत “कंटाळवाणी तपासणी” स्पष्टपणे दिली आहे.}

\emph{त्या} कष्टाच्या सरावाकडे वळण्यापूर्वी आणखी काही ``तात्काळ'' गोष्टी
तपासल्या पाहिजेत. उदाहरणार्थ, कोणत्याही $\alpha$ आणि $\beta$ छेदांसाठी
व्याख्यायित केलेला $\alpha + \beta$ खरोखरच \emph{छेद} आहे याची खात्री
हवी. त्याची सिद्धता अशी:

\begin{proof}	
$\alpha$ आणि $\beta$ हे दोन्ही छेद असल्यामुळे $\alpha + \beta = \Setabs{p
+ q}{p \in \alpha \land q \in \beta}$ हा $\Rat$ चा अरिक्त उचित उपसंच
आहे. आता काही $p \in \alpha$ आणि $q \in \beta$ साठी $x < p + q$ असे
समजा. मग $x - p < q$, म्हणून $x - p \in \beta$, आणि $x = p + (x - p)
\in \alpha + \beta$. म्हणून $\alpha + \beta$ हा $\Rat$ चा आरंभीचा खंड
आहे. शेवटी, कोणत्याही $p + q \in \alpha + \beta$ साठी, $\alpha$ आणि
$\beta$ हे दोन्ही छेद असल्यामुळे $p < p_1$ आणि $q < q_1$ असे काही
$p_1 \in \alpha$ आणि $q_1 \in \beta$ असतात; म्हणून $p + q < p_1 + q_1
\in \alpha + \beta$; त्यामुळे $\alpha + \beta$ ला महत्तम घटक नाही.
\end{proof}

अशाच प्रयत्नांनी $\alpha - \beta$, $\alpha \times \beta$ आणि $\alpha
\div \beta$ हे छेद आहेत हे तुम्हाला तपासता येईल (शेवटच्या प्रकरणात
$\beta$ हा शून्य-छेद असलेले प्रकरण वगळून). पुन्हा, हे काम आम्ही
तुमच्यासाठीच सोडतो.
%READERNOTE{MRARITH-008: गोठवलेल्या इंग्रजी स्रोताच्या आधीच्या cuts विभागात वजाबाकी आणि भागाकाराची सूत्रे निष्क्रिय टिप्पण्यांतच आहेत; त्यामुळे हे वाक्य सक्रियपणे पूर्ण व्याख्या न दिलेल्या क्रियांना गृहीत धरते. मराठी मजकूर दावा जतन करतो आणि ही स्रोत-अपूर्णता स्पष्ट करतो.}

\begin{prob}
$\Real$ हे क्रमित क्षेत्र आहे हे सिद्ध करा.
\end{prob}

पण एक लहानसा अपूर्ण मुद्दा आवरायचा आहे. \olref[cuts]{sec} मध्ये आपण
$\sqrt{2} = \Setabs{p \in \Rat}{p < 0 \text{ किंवा }p^2 < 2}$ असे घेता
येईल असे सुचवले. मात्र हा संच \emph{छेद} आहे हे दाखवणे आवश्यक आहे.
त्याची सिद्धता अशी:

\begin{proof}
हा परिमेय संख्यांचा अरिक्त उचित आरंभीचा खंड आहे हे स्पष्ट आहे; म्हणून
त्याला महत्तम घटक नाही एवढे दाखवणे पुरेसे आहे. विशेषतः, $p$ ही $p^2 < 2$
असलेली धन परिमेय संख्या आणि $q = \frac{2p+2}{p+2}$ असेल, तर $p < q$
आणि $q^2 < 2$ हे दोन्ही दाखवणे पुरेसे आहे. $p < q$ हे पाहण्यासाठी फक्त
पुढील नोंद घ्या:
\begin{align*}
	p^2 &< 2\\
	p^2 + 2p &< 2 + 2p\\
	p(p + 2) &< 2 + 2p\\
	p &< \tfrac{2+2p}{p+2} = q
\end{align*}
$q^2 < 2$ हे पाहण्यासाठी फक्त पुढील नोंद घ्या:
\begin{align*}
	p^2 &< 2\\
	2p^2 + 4p + 2 &< p^2 + 4p+ 4\\
	4p^2 + 8p + 4 &< 2(p^2 + 4p + 4)\\
	(2p+2)^2 & <2(p+2)^2\\
	\tfrac{(2p+2)^2}{(p+2)^2} &< 2\\
	q^2 &< 2
\end{align*}
\end{proof}
\end{document}
